\begin{frame}{자주 묻는 질문(15.1)}
  \small
  \begin{itemize}
    \item \textbf{Q: Week 4 가치반복도 ``모델기반''인데, 왜 이 장에서
      다시 이 용어?} --- 차이는 \textbf{스케일과 쓰임새}: Week 4는
      상태가 적어 정책 \textbf{전체}를 미리 계산, 이번 장(Dyna-Q,
      MCTS)은 상태가 너무 많아 ``지금 이 상태에서 무엇을 할지''
      \textbf{만} 계산
    \item \textbf{Q: 모델기반이 model-free보다 항상 좋은가?} ---
      \textbf{아니} --- 정확한(또는 잘 학습된) 모델이 있을 때 데이터
      효율이 훨씬 좋지만, 모델 학습 비용과 \textbf{오차 누적} 위험이
      있음. 로봇 팔처럼 물리 법칙이 복잡한 환경은 Week 9--13의
      model-free 쪽이 더 실용적인 경우 많음
    \item \textbf{Q: known model이면 learned model 얘기를 안 해도
      되나?} --- 실제 환경 대부분이 바둑처럼 규칙이 완전히 알려진
      게임이 \textbf{아님}(로봇·화학 공정·교통·인간-기계 상호작용).
      ``공짜 모델''은 예외적 특권, \textbf{일반은 learned model} ---
      2010년대 이후 연구(world model 계열)의 주목적도 이 케이스
      (15.4절)
  \end{itemize}
\end{frame}
